\chapter{Bilan et décision}

\section{Bilan de la campagne technique}

\begin{table}[H]
\centering
\caption{État de validation au 17/08/2026}
\begin{tabularx}{\textwidth}{|L{4.4cm}|C{3cm}|Y|}
\hline
\rowcolor{ctmint}\textbf{Axe} & \textbf{État} & \textbf{Commentaire} \\ \hline
Backend et règles métier & \statuspass & 270 tests et 2\,784 assertions. \\ \hline
Gateway OCPP & \statuspass & 23 tests exécutés dans l'image Docker dédiée. \\ \hline
Outils simulateur & \statuspass & 9 tests Node réussis. \\ \hline
Frontend & \statuspass & 46 tests, lint et build TypeScript/Vite réussis. \\ \hline
Infrastructure et CI/CD & \statuspass & 30 tests de configuration, sécurité, Docker et Azure réussis. \\ \hline
Sécurité des dépendances PHP & \statuspass & Aucun avis Composer connu après mise à jour ciblée. \\ \hline
Recette fonctionnelle & \statusmanual & Les fiches du présent document doivent être renseignées. \\ \hline
Performance et charge & \statusmanual & Seuils et volumes à valider avec l'encadrant. \\ \hline
Borne physique et paiement réel & \statusdeferred & Hors MVP simulé actuel. \\ \hline
Firmware & \statusdeferred & US-38 reportée. \\ \hline
\end{tabularx}
\end{table}

\section{Risques résiduels}

\begin{enumerate}[leftmargin=7mm]
  \item les tests frontend couvrent les contrats critiques, mais les parcours
        E2E visuels multi-navigateurs restent à exécuter manuellement ;
  \item le worker PDF et certains modules spécialisés restent volumineux ;
  \item les seuils Heartbeat, SLA et disponibilité doivent être recalibrés avec
        l'usage réel ;
  \item les comportements d'un prestataire de paiement et d'une borne physique
        ne sont pas couverts par les simulateurs ;
  \item les tests de charge, la restauration complète et la supervision longue
        durée restent nécessaires malgré le déploiement HTTPS actuel.
\end{enumerate}

\section{Décision proposée}

\begin{testbox}{Décision technique}
La base technique du MVP est \textbf{apte à entrer en recette fonctionnelle}.
Cette décision ne constitue pas encore une acceptation finale : le
procès-verbal de recette devra reprendre les cas manuels critiques, les
anomalies éventuelles, les réserves et les signatures des parties.
\end{testbox}

\section{Étapes suivantes}

\begin{enumerate}[leftmargin=7mm]
  \item exécuter d'abord les cas CT-AUTH-03, CT-STN-01, CT-OCPP-01,
        CT-SIM-01, CT-CHG-01, CT-PAY-01, CT-OPS-01, CT-COM-01,
        CT-SEC-01 et CT-DEP-01 ;
  \item corriger toute anomalie S1/S2 puis relancer les suites associées ;
  \item compléter la campagne responsive, multi-navigateurs et performance ;
  \item produire le procès-verbal de recette à partir des résultats signés ;
  \item mettre à jour ce cahier à chaque changement du backlog ou du contrat API.
\end{enumerate}
